\section{Endereços, entropia, descrição e busca}
\subsection{Comprimento do índice canônico}
Para um átomo com parâmetros $(k,m,y_0,s)$, defina
\eqn{46}{H=\max\{2,k+1,|m|+1,\heightq(y_0),\heightq(s)\}.}
Escreva $\LL_H\subset\LL$ para o conjunto dos átomos de altura no máximo $H$.

\begin{teorema}[cota polinomial do endereço]
O único endereço $n$ do átomo satisfaz
\eqn{47}{n=O(H^8),\qquad\log_2(n+1)=O(\log H).}
As constantes são universais para a enumeração fixada.
\end{teorema}
\begin{proof}
Pelo lema~2.4, os postos $r$ de $y_0$ e $t$ de $s$ são $O(H^2)$. Pelo lema~2.2, $\mu=Z^{-1}(m)=O(H)$, e $k=O(H)$. A cota do lema~2.1 dá
\[
u=\pi(k,\mu)=O(H^2),\qquad v=\pi(r,t)=O(H^4),
\]
e, por uma última aplicação,
\[
n=\pi(u,v)=O\bigl((u+v+1)^2\bigr)=O(H^8).
\]
Como $H\ge2$, tomar logaritmos fornece a segunda afirmação.
\end{proof}

\subsection{Contagem por altura e consequência minimax para o endereço}
A contagem dos átomos determina o menor pior índice de uma enumeração; o aninhamento de Cantor determina o crescimento específico de $\Gamma$.
\begin{lema}[cardinalidade das cascas de altura]
Para todo $H\ge2$,
\eqn{48}{|\LL_H|\le18H^6.}
Além disso, existem uma constante universal $c_{\LL}>0$ e $H_1\ge2$ tais que, para todo $H\ge H_1$,
\eqn{49}{|\LL_H|\ge c_{\LL}H^6.}
Em particular, $|\LL_H|=\Theta(H^6)$.
\end{lema}
\begin{proof}
Pelo teorema~3.6, um átomo é determinado bijetivamente pela quádrupla $(k,m,y_0,s)\in\N_0\times\Z\times\Q\times\Q^*$, e a altura \eqref{eq:46} limita cada coordenada separadamente. Assim
\[
|\LL_H|=\#\{k\le H-1\}\cdot\#\{|m|\le H-1\}\cdot|\Q_H|\cdot|\Q_H\setminus\{0\}|,
\]
com $\Q_H$ como no lema~2.3. O primeiro fator é $H$, o segundo é $2H-1$, e os dois últimos estão entre $c_1H^2$ e $3H^2$ para $H$ grande. Logo
\begin{align*}
|\LL_H|&\le H\cdot2H\cdot3H^2\cdot3H^2=18H^6,\\
|\LL_H|&\ge H\cdot H\cdot c_1H^2\cdot(c_1H^2-1)\ge\frac{c_1^2}2H^6
\end{align*}
para todo $H$ suficientemente grande. Tome $c_{\LL}=c_1^2/2$.
\end{proof}

Com orçamentos desacoplados,
\[
\LL_{K_0,M_0,Y,S}:=\{(k,m,y_0,s):k\le K_0,\ |m|\le M_0,\ \heightq(y_0)\le Y,\ \heightq(s)\le S\},
\]
a bijeção do teorema~3.6 dá
\[
|\LL_{K_0,M_0,Y,S}|=(K_0+1)(2M_0+1)|\Q_Y|(|\Q_S|-1).
\]
A ordem é $(K_0+1)(M_0+1)Y^2S^2$ para $Y,S$ grandes. O calibre \eqref{eq:46} acopla $K_0=M_0=H-1$ e $Y=S=H$, produzindo $H^6$; já $K_0=M_0=\lfloor\log_2H\rfloor$ e $Y=S=H$ produziriam $H^4(\log H)^2$.

\begin{proposicao}[ordem exata do pior endereço aritmético]
Se
\[
M_\Gamma(H):=\max_{a\in\LL_H}\Gamma^{-1}(a),
\]
então
\eqn{50}{M_\Gamma(H)=\Theta(H^8).}
Consequentemente, o comprimento da maior palavra do gerador aritmético sobre $\LL_H$ é $8\log_2H+O(1)$ bits.
\end{proposicao}
\begin{proof}
A cota superior é o teorema~7.1. Para a inferior, escreva $N_H=|\Q_H|$. Pelo lema~2.3, $N_H\ge c_1H^2$ para $H$ suficientemente grande. Como as enumerações de \eqref{eq:05} listam por altura, existem $y_H\in\Q_H$ e $s_H\in\Q_H\setminus\{0\}$ cujos postos respectivos satisfazem
\[
r_H=N_H-1,\qquad t_H=N_H-2,
\]
e, portanto, $r_H,t_H\ge c_2H^2$ para alguma constante universal $c_2>0$. Considere o átomo de parâmetros $(k,m,y_0,s)=(0,0,y_H,s_H)$, que pertence a $\LL_H$. Como
\[
\pi(x,y)=\frac{(x+y)(x+y+1)}2+y\ge\frac{(x+y)^2}2,
\]
temos $v_H=\pi(r_H,t_H)\ge c_3H^4$. Além disso, $\pi(0,0)=0$, de modo que o endereço desse átomo satisfaz
\[
n_H=\pi\bigl(\pi(0,0),v_H\bigr)=\pi(0,v_H)\ge\frac{v_H^2}2\ge c_4H^8.
\]
Logo $M_\Gamma(H)=\Omega(H^8)$, que, junto da cota superior, prova \eqref{eq:50}. A afirmação em bits segue tomando logaritmos.
\end{proof}

\begin{teorema}[cota minimax relativa ao calibre]
Seja $A:\LL\to\N_0$ uma aplicação injetiva qualquer. Para todo $H\ge2$,
\eqn{51}{\max_{a\in\LL_H}A(a)\ge|\LL_H|-1.}
Consequentemente, para todo $H\ge H_1$, existe um átomo de altura no máximo $H$ com
\eqn{52}{A(\text{átomo})\ge c_{\LL}H^6-1,\qquad
\log_2\bigl(A(\text{átomo})+1\bigr)\ge6\log_2H-O(1).}
Nenhum endereçamento injetivo usa $o(\log H)$ bits no pior caso sobre $\LL_H$.
\end{teorema}
\begin{proof}
Por injetividade, $A(\LL_H)$ é um subconjunto de $\N_0$ com $|\LL_H|$ elementos distintos; seu máximo é, portanto, pelo menos $|\LL_H|-1$, o que prova \eqref{eq:51}. Aplique então \eqref{eq:49} e tome logaritmos.
\end{proof}
\begin{observacao}
O argumento vale para todo endereçamento injetivo, inclusive os definidos por tabelas, estado ou escolhas arbitrárias. Com a proposição~7.3, ele compara $8\log_2H+O(1)$ bits de $\Gamma$ a $6\log_2H+O(1)$ bits do limite minimax: a razão assintótica dos comprimentos é $4/3$.
\end{observacao}

\begin{definicao}[enumeração por altura da biblioteca]
Para $H\ge2$, a casca $\LL_H\setminus\LL_{H-1}$ é finita e não vazia. Ordene $\LL$ listando as cascas em ordem crescente de $H$ e, dentro de cada casca, ordenando as quádruplas $(k,m,y_0,s)$ lexicograficamente por
\[
\bigl(k,Z^{-1}(m),\rank_Q(y_0),\rank_{Q^*}(s)\bigr),
\]
onde $\rank$ denota o posto nas enumerações \eqref{eq:05}. A bijeção resultante é denotada $\Gamma^\sharp:\N_0\to\LL$, e $n^\sharp$ é o endereço de um átomo nessa ordem.
\end{definicao}

\begin{teorema}[endereço minimax exato para a exaustão]
A aplicação $\Gamma^\sharp$ é uma bijeção de $\N_0$ sobre $\LL$, computável nos dois sentidos, e para todo $H\ge2$
\eqn{53}{\max_{a\in\LL_H}n^\sharp(a)=|\LL_H|-1\le18H^6-1.}
Além disso, para qualquer endereçamento injetivo $A:\LL\to\N_0$,
\eqn{54}{\max_{a\in\LL_H}A(a)\ge|\LL_H|-1.}
Logo $\Gamma^\sharp$ realiza o mínimo possível do maior índice sobre $\LL_H$. Esse pior índice é $\Theta(H^6)$ e seu comprimento binário é minimax-ótimo para o calibre \eqref{eq:46}.
\end{teorema}
\begin{proof}
Cada casca é finita pelo lema~7.2 e a ordem dentro dela é total, de modo que a concatenação está bem definida e é bijetiva; a computabilidade nos dois sentidos segue de as cascas serem finitas e explicitamente enumeráveis. Por construção, $\Gamma^\sharp$ lista primeiro, e exatamente, os elementos de $\LL_H$; portanto $\Gamma^\sharp$ restrita a $\{0,\ldots,|\LL_H|-1\}$ é uma bijeção sobre $\LL_H$, o que dá a igualdade em \eqref{eq:53}, e \eqref{eq:48} dá a cota numérica. A desigualdade \eqref{eq:54} é \eqref{eq:51}: a imagem $A(\LL_H)$ tem $|\LL_H|$ elementos distintos de $\N_0$, logo seu máximo é ao menos $|\LL_H|-1$.
\end{proof}

\begin{observacao}[endereçamento e armazenamento]
$\Gamma$ descompacta o índice em célula, valor ancorado e inclinação pelo pareamento de Cantor. $\Gamma^\sharp$ ordena as mesmas quádruplas por altura e atinge o menor pior índice. A primeira usa três raízes quadradas inteiras na descompactação de \eqref{eq:13}; a segunda calcula cardinalidades de cascas, com somas de totientes truncadas e custo polinomial em $H$.

A conversão entre as duas ordens decodifica a quádrupla e a recodifica, preservando a célula e a função afim. A biblioteca tem, assim, a semântica fixada por $\Gamma$ e um codec minimax dado por $\Gamma^\sharp$.
\end{observacao}
Sob domínio, raiz e convenção de fronteira compartilhados, os endereços recuperam as células. O conjunto das folhas determina a árvore podada por fechamento ancestral até a raiz, e seus extremos determinam a ordem espacial. A transmissão usa apenas $N$, $H$ e a lista de endereços.

\begin{teorema}[serialização condicionada de montagens]
Fixe, como informação compartilhada, o domínio $K$, a raiz diádica inicial, a convenção de fronteira e a enumeração usada. Seja $g\in\Asm(K)$ descrita por uma árvore diádica com $N$ folhas, e seja $H\ge2$ a maior altura de seus átomos. Existe um fluxo de bits autodelimitado, decodificável sem conhecer previamente $N$ ou $H$, com comprimento
\eqn{55}{\begin{aligned}
L(g\mid K)&\le(2\lfloor\log_2N\rfloor+1)+(2\lfloor\log_2H\rfloor+1)
+N\left\lceil\log_2(18H^6)\right\rceil\\
&=N\bigl(6\log_2H+O(1)\bigr)+O(\log N+\log H)
\end{aligned}}
quando os átomos são transmitidos por $\Gamma^\sharp$. Com o gerador aritmético $\Gamma$, o termo por folha é $8\log_2H+O(1)$ bits.
\end{teorema}
\begin{proof}
Codifique $N$ e $H$, nessa ordem, pelo código gama de Elias, com comprimentos $2\lfloor\log_2N\rfloor+1$ e $2\lfloor\log_2H\rfloor+1$. Conhecido $H$, cada endereço de $\Gamma^\sharp$ ocupa exatamente $b_H=\lceil\log_2(18H^6)\rceil$ bits, pois $n^\sharp<18H^6$ pelo teorema~7.7. O cabeçalho $N$ informa quantos blocos de largura $b_H$ devem ser lidos. Os endereços recuperam os átomos e suas células; ordenar os extremos fornece a ordem espacial, e fechar as células por ancestrais recupera a árvore podada. O domínio compartilhado determina apenas o recorte nas extremidades. Isso prova a primeira cota e a decodificação sem esqueleto separado. Para $\Gamma$, use $\max_{a\in\LL_H}\Gamma^{-1}(a)=\Theta(H^8)$ da proposição~7.3.
\end{proof}
\begin{observacao}
Para extremos reais arbitrários, o domínio $K$ pode exigir informação infinita. Por isso, o teorema~7.9 mede a descrição condicionada ao domínio e à convenção de fronteira declarados.
\end{observacao}

\subsection{Três formatos prefix-free e propagação de altura}
Para $n\ge1$, seja
\[
\gamma(n)=0^{\lfloor\log_2n\rfloor}\operatorname{bin}(n)
\]
o código gama de Elias, e ponha
\[
\zeta(z)=\begin{cases}2z,&z\ge0,\\-2z-1,&z<0.\end{cases}
\]
Para $r=p/q$ reduzido, $q>0$, e para o diádico escrito canonicamente como $t=u2^{-k}$, com $u$ ímpar quando $k>0$, defina
\eqn{56}{C_{\Z}(z)=\gamma(\zeta(z)+1),}
\eqn{57}{C_{\rm rat}(p/q)=C_{\Z}(p)\gamma(q),\qquad L_{\rm rat}(p/q)=|C_{\rm rat}(p/q)|,}
\eqn{58}{C_{\rm dyad}(u,k)=C_{\Z}(u)\gamma(k+1),\qquad L_{\rm dyad}(u,k)=|C_{\rm dyad}(u,k)|.}
Cada componente é prefix-free e o cabeçalho determina a quantidade de componentes, tornando as concatenações decodificáveis sem separadores.

\begin{lema}[cálculo elementar de alturas]
Para quaisquer $r,s\in\Q$,
\eqn{59}{\heightq(r\pm s)\le2\heightq(r)\heightq(s),\qquad\heightq(rs)\le\heightq(r)\heightq(s).}
Em particular, se ambos têm altura no máximo $H$, as cotas respectivas são $2H^2$ e $H^2$. Para uma montagem de altura atômica máxima $H$, todo nó interno $t_j=u_j2^{-k_j}$ satisfaz $|u_j|\le H$ e $k_j\le H$ em sua forma reduzida; além disso,
\eqn{60}{\heightq(s_j-s_{j-1})\le2H^2,}
\eqn{61}{\heightq(s_jt_j)\le H2^H,}
\eqn{62}{\heightq\bigl(g(t_0)-s_0t_0\bigr)\le2^{H+1}H^2.}
\end{lema}
\begin{proof}
Escreva $r=p/q$ e $s=a/b$ em forma reduzida e ponha $R=\heightq(r)$, $S=\heightq(s)$. Antes de qualquer cancelamento, $r\pm s=(pb\pm aq)/(qb)$ e $rs=pa/(qb)$; os numeradores têm módulos no máximo $2RS$ e $RS$, respectivamente, e os denominadores são no máximo $RS$. Reduzir não aumenta altura, o que prova \eqref{eq:59}. Um nó é fronteira de uma célula $J_{k,m}$, logo sua forma não reduzida tem expoente $k\le H-1$ e numerador $m$ ou $m+1$, de módulo no máximo $H$; redução só diminui esses dados. As cotas \eqref{eq:60} e \eqref{eq:61} seguem da primeira parte e de $\heightq(t_j)\le2^H$. Finalmente, na primeira folha, $g(t_0)-s_0t_0=y_0-s_0x_0$; como $\heightq(x_0)\le2^H$, duas aplicações das cotas de produto e soma dão \eqref{eq:62}.
\end{proof}
Fixe a ordem espacial das folhas. Os três fluxos executáveis são:
\begin{enumerate}[label=(\roman*),topsep=5pt,itemsep=4pt]
\item $C_{\rm atom}(g)$: $\gamma(N)$, seguido, para cada folha, de $\gamma(k+1)C_{\Z}(m)C_{\rm rat}(y_0)C_{\rm rat}(s)$;
\item $C_{\rm ReLU,str}(g)$: $\gamma(N)$, o primeiro átomo na forma ancorada $C_{\rm dyad}(x_0)C_{\rm rat}(y_0)C_{\rm rat}(s_0)$, e os pares $C_{\rm dyad}(t_j)C_{\rm rat}(s_j-s_{j-1})$;
\item $C_{\rm ReLU,rat}(g)$: $\gamma(N)$ e a forma expandida $C_{\rm rat}(c)C_{\rm rat}(s_0)$, seguida dos pares $C_{\rm rat}(t_j)C_{\rm rat}(s_j-s_{j-1})$, onde $c=g(t_0)-s_0t_0$.
\end{enumerate}
Denote os comprimentos por $L_{\rm atom}$, $L_{\rm ReLU,str}$ e $L_{\rm ReLU,rat}$. O primeiro fluxo é uma alternativa estrutural à serialização por $\Gamma^\sharp$ do teorema~7.9; ambos têm a mesma ordem de comprimento.

\begin{teorema}[comparação rigorosa dos formatos]
Os três fluxos são prefix-free e possuem decodificadores exatos. Para toda montagem com $N$ folhas e altura atômica máxima $H\ge2$,
\eqn{63}{L_{\rm atom}(g)=O(N\log H),\quad L_{\rm ReLU,str}(g)=O(N\log H),\quad
L_{\rm ReLU,rat}(g)=O(NH),}
com constantes universais. A última ordem é atingida conjuntamente em $N$ e $H$: para todo $H\ge2$ e todo $N$ par com $2\le N\le H$, considere as folhas $J_{H-1,j}$, $0\le j<N$, com valores centrais zero e inclinações $(-1)^j$. Essa montagem contínua tem altura $H$, e
\eqn{64}{L_{\rm atom}=O(N\log H),\quad L_{\rm ReLU,str}=O(N\log H),\quad
L_{\rm ReLU,rat}=\Omega(NH).}
A equivalência funcional montagem--ReLU preserva, portanto, a ordem de bits no formato estrutural fixado.
\end{teorema}
\begin{proof}
O código gama é prefix-free; o mapa $\zeta$ é bijetivo entre $\Z$ e $\N_0$, e frações reduzidas e pares diádicos canônicos são únicos. Lendo o cabeçalho $N$, cada decodificador sabe quantos blocos seguintes consumir. No fluxo atômico, os blocos recuperam exatamente $(k,m,y_0,s)$ de cada folha. No fluxo ReLU estrutural, o primeiro bloco recupera $(x_0,y_0,s_0)$ e, portanto, $c=y_0-s_0x_0$; os demais recuperam os nós e saltos de inclinação da equação \eqref{eq:24}; o terceiro fluxo já transmite esses parâmetros expandidos. Logo os decodificadores são exatos.

Se $|z|\le A$, então $|C_{\Z}(z)|=O(\log(1+A))$, e se $\heightq(r)\le A$, então $L_{\rm rat}(r)=O(\log(1+A))$. As coordenadas dos átomos têm tamanho $O(H)$; os pares diádicos estruturais têm numerador e expoente $O(H)$; e os saltos têm altura $O(H^2)$ pelo lema~7.11. Isso prova as duas primeiras cotas de \eqref{eq:63}. Na forma ingênua, cada nó pode ter denominador $2^H$, e o intercepto pode ter altura $2^{H+1}H^2$; o lema~7.11 dá $O(H)$ bits por bloco e prova a terceira cota.

Na família final, as duas peças adjacentes ao nó $t_j=j2^{-(H-1)}$ assumem ali o mesmo valor $(-1)^{j-1}2^{-H}$, logo a montagem é contínua. Como $j<N\le H$, sua altura atômica máxima é exatamente $H$. Há ao menos $(N-2)/2$ índices internos ímpares $j$; para eles, $t_j$ já está reduzido e tem denominador $2^{H-1}$. Cada bloco racional correspondente custa $\Omega(H)$ bits por \eqref{eq:57}, e os saltos de inclinação são não nulos. Portanto $L_{\rm ReLU,rat}=\Omega(NH)$; os blocos atômicos e os pares diádicos estruturais custam $O(\log H)$ cada, dando as outras duas cotas. Para $N=2$, o intercepto $c=-2^{-H}$ fornece diretamente a mesma cota $\Omega(H)$. Isso prova \eqref{eq:64}.
\end{proof}
\input{secoes/07b_erro_entropia}
\input{secoes/07c_catalogos}
